% Options for packages loaded elsewhere
\PassOptionsToPackage{unicode}{hyperref}
\PassOptionsToPackage{hyphens}{url}
\PassOptionsToPackage{dvipsnames,svgnames,x11names}{xcolor}
%
\documentclass[
  10pt,
]{article}
\usepackage{amsmath,amssymb}
\usepackage{iftex}
\ifPDFTeX
  \usepackage[T1]{fontenc}
  \usepackage[utf8]{inputenc}
  \usepackage{textcomp} % provide euro and other symbols
\else % if luatex or xetex
  \usepackage{unicode-math} % this also loads fontspec
  \defaultfontfeatures{Scale=MatchLowercase}
  \defaultfontfeatures[\rmfamily]{Ligatures=TeX,Scale=1}
\fi
\usepackage{lmodern}
\ifPDFTeX\else
  % xetex/luatex font selection
  \setmainfont[]{DejaVu Serif}
  \setsansfont[]{DejaVu Sans}
  \setmonofont[]{DejaVu Sans Mono}
  \setmathfont[]{Latin Modern Math}
\fi
% Use upquote if available, for straight quotes in verbatim environments
\IfFileExists{upquote.sty}{\usepackage{upquote}}{}
\IfFileExists{microtype.sty}{% use microtype if available
  \usepackage[]{microtype}
  \UseMicrotypeSet[protrusion]{basicmath} % disable protrusion for tt fonts
}{}
\makeatletter
\@ifundefined{KOMAClassName}{% if non-KOMA class
  \IfFileExists{parskip.sty}{%
    \usepackage{parskip}
  }{% else
    \setlength{\parindent}{0pt}
    \setlength{\parskip}{6pt plus 2pt minus 1pt}}
}{% if KOMA class
  \KOMAoptions{parskip=half}}
\makeatother
\usepackage{xcolor}
\usepackage[margin=2.0cm]{geometry}
\usepackage{color}
\usepackage{fancyvrb}
\newcommand{\VerbBar}{|}
\newcommand{\VERB}{\Verb[commandchars=\\\{\}]}
\DefineVerbatimEnvironment{Highlighting}{Verbatim}{commandchars=\\\{\}}
% Add ',fontsize=\small' for more characters per line
\newenvironment{Shaded}{}{}
\newcommand{\AlertTok}[1]{\textcolor[rgb]{1.00,0.00,0.00}{\textbf{#1}}}
\newcommand{\AnnotationTok}[1]{\textcolor[rgb]{0.38,0.63,0.69}{\textbf{\textit{#1}}}}
\newcommand{\AttributeTok}[1]{\textcolor[rgb]{0.49,0.56,0.16}{#1}}
\newcommand{\BaseNTok}[1]{\textcolor[rgb]{0.25,0.63,0.44}{#1}}
\newcommand{\BuiltInTok}[1]{\textcolor[rgb]{0.00,0.50,0.00}{#1}}
\newcommand{\CharTok}[1]{\textcolor[rgb]{0.25,0.44,0.63}{#1}}
\newcommand{\CommentTok}[1]{\textcolor[rgb]{0.38,0.63,0.69}{\textit{#1}}}
\newcommand{\CommentVarTok}[1]{\textcolor[rgb]{0.38,0.63,0.69}{\textbf{\textit{#1}}}}
\newcommand{\ConstantTok}[1]{\textcolor[rgb]{0.53,0.00,0.00}{#1}}
\newcommand{\ControlFlowTok}[1]{\textcolor[rgb]{0.00,0.44,0.13}{\textbf{#1}}}
\newcommand{\DataTypeTok}[1]{\textcolor[rgb]{0.56,0.13,0.00}{#1}}
\newcommand{\DecValTok}[1]{\textcolor[rgb]{0.25,0.63,0.44}{#1}}
\newcommand{\DocumentationTok}[1]{\textcolor[rgb]{0.73,0.13,0.13}{\textit{#1}}}
\newcommand{\ErrorTok}[1]{\textcolor[rgb]{1.00,0.00,0.00}{\textbf{#1}}}
\newcommand{\ExtensionTok}[1]{#1}
\newcommand{\FloatTok}[1]{\textcolor[rgb]{0.25,0.63,0.44}{#1}}
\newcommand{\FunctionTok}[1]{\textcolor[rgb]{0.02,0.16,0.49}{#1}}
\newcommand{\ImportTok}[1]{\textcolor[rgb]{0.00,0.50,0.00}{\textbf{#1}}}
\newcommand{\InformationTok}[1]{\textcolor[rgb]{0.38,0.63,0.69}{\textbf{\textit{#1}}}}
\newcommand{\KeywordTok}[1]{\textcolor[rgb]{0.00,0.44,0.13}{\textbf{#1}}}
\newcommand{\NormalTok}[1]{#1}
\newcommand{\OperatorTok}[1]{\textcolor[rgb]{0.40,0.40,0.40}{#1}}
\newcommand{\OtherTok}[1]{\textcolor[rgb]{0.00,0.44,0.13}{#1}}
\newcommand{\PreprocessorTok}[1]{\textcolor[rgb]{0.74,0.48,0.00}{#1}}
\newcommand{\RegionMarkerTok}[1]{#1}
\newcommand{\SpecialCharTok}[1]{\textcolor[rgb]{0.25,0.44,0.63}{#1}}
\newcommand{\SpecialStringTok}[1]{\textcolor[rgb]{0.73,0.40,0.53}{#1}}
\newcommand{\StringTok}[1]{\textcolor[rgb]{0.25,0.44,0.63}{#1}}
\newcommand{\VariableTok}[1]{\textcolor[rgb]{0.10,0.09,0.49}{#1}}
\newcommand{\VerbatimStringTok}[1]{\textcolor[rgb]{0.25,0.44,0.63}{#1}}
\newcommand{\WarningTok}[1]{\textcolor[rgb]{0.38,0.63,0.69}{\textbf{\textit{#1}}}}
\usepackage{longtable,booktabs,array}
\usepackage{calc} % for calculating minipage widths
% Correct order of tables after \paragraph or \subparagraph
\usepackage{etoolbox}
\makeatletter
\patchcmd\longtable{\par}{\if@noskipsec\mbox{}\fi\par}{}{}
\makeatother
% Allow footnotes in longtable head/foot
\IfFileExists{footnotehyper.sty}{\usepackage{footnotehyper}}{\usepackage{footnote}}
\makesavenoteenv{longtable}
\setlength{\emergencystretch}{3em} % prevent overfull lines
\providecommand{\tightlist}{%
  \setlength{\itemsep}{0pt}\setlength{\parskip}{0pt}}
\setcounter{secnumdepth}{-\maxdimen} % remove section numbering
% doc_header.tex -- LaTeX preamble additions used by build_docs.sh for every
% Markdown -> LaTeX -> PDF conversion (supplements, README, submission files).
\usepackage{etoolbox}
% longtable may not be loaded by pandoc when a document has no table
\usepackage{longtable}
% bold vectors: \bm -> \symbf under unicode-math (XeLaTeX), bm package otherwise
\makeatletter
\@ifpackageloaded{unicode-math}{\providecommand{\bm}[1]{\symbf{#1}}}{\usepackage{bm}}
\makeatother
% smaller type in long tables so that wide audit tables fit the text width;
% inside a table a line break is allowed after an underscore of plain text
% (pandoc writes it as \_), so that long identifiers wrap inside their column
\AtBeginEnvironment{longtable}{\footnotesize
  \let\sptunderscore\_\def\_{\sptunderscore\allowbreak}}
\setlength{\LTcapwidth}{\textwidth}
% allow line breaks after slashes and hyphens inside table cells
\usepackage{xurl}
% code blocks: break over-long lines instead of running into the margin.
% Plain ``` blocks are typeset by pandoc as {verbatim}, highlighted ones
% (```bash) as {Highlighting}; both are re-declared as breakable fancyvrb
% environments.  The continuation line is marked with a small hooked arrow.
\usepackage{fvextra}
\RecustomVerbatimEnvironment{verbatim}{Verbatim}{breaklines,breakanywhere,fontsize=\small}
\makeatletter
\@ifundefined{Highlighting}{}{%
  \RecustomVerbatimEnvironment{Highlighting}{Verbatim}{breaklines,breakanywhere,fontsize=\small,commandchars=\\\{\}}}
\makeatother
% inline code (pandoc writes it as \texttt{...}): never hyphenate -- an inserted
% hyphen could be read as part of a file name -- but allow a line break after
% an underscore, slash, comma, colon or equals sign, so that long paths and
% column lists wrap inside the text width and inside table columns.
\ExplSyntaxOn
\tl_new:N \l__spt_code_tl
\RenewDocumentCommand \texttt { m }
  {
    \tl_set:Nn \l__spt_code_tl {#1}
    \regex_replace_all:nnN { ([\/\,\:\=]) } { \1 \c{allowbreak} } \l__spt_code_tl
    \regex_replace_all:nnN { \c{_} } { \c{_} \c{allowbreak} } \l__spt_code_tl
    \group_begin: \ttfamily \hyphenchar \font = -1 \scan_stop: \l__spt_code_tl \group_end:
  }
\ExplSyntaxOff
% A figure is written in the Markdown source as an image paragraph followed by
% a caption paragraph ("**Figure S1.** ..."), so that the caption is also
% visible on GitHub.  An image that starts a paragraph is centered.
\usepackage{graphicx}% pandoc loads it only for documents that contain an image
\NewCommandCopy\sptincludegraphics\includegraphics
\RenewDocumentCommand\includegraphics{O{} m}{%
  \ifvmode\begin{center}\sptincludegraphics[#1]{#2}\end{center}%
  \else\sptincludegraphics[#1]{#2}\fi}
% keep figures where they are referenced
\usepackage{float}
\floatplacement{figure}{H}
% running header with the document title
\usepackage{fancyhdr}
\pagestyle{fancy}
\fancyhf{}
\fancyhead[L]{\small PROJECT spacetimes --- documentation and provenance}
\fancyhead[R]{\small\thepage}
\renewcommand{\headrulewidth}{0.4pt}
\ifLuaTeX
  \usepackage{selnolig}  % disable illegal ligatures
\fi
\IfFileExists{bookmark.sty}{\usepackage{bookmark}}{\usepackage{hyperref}}
\IfFileExists{xurl.sty}{\usepackage{xurl}}{} % add URL line breaks if available
\urlstyle{same}
\hypersetup{
  pdftitle={P0 --- Why this project is not pushed by the assistant, and how it reaches GitHub},
  pdfauthor={Prepared for Patrick Lee Nash, Ph.D.},
  colorlinks=true,
  linkcolor={blue},
  filecolor={Maroon},
  citecolor={blue},
  urlcolor={blue},
  pdfcreator={LaTeX via pandoc}}

\title{P0 --- Why this project is not pushed by the assistant, and how
it reaches GitHub}
\usepackage{etoolbox}
\makeatletter
\providecommand{\subtitle}[1]{% add subtitle to \maketitle
  \apptocmd{\@title}{\par {\large #1 \par}}{}{}
}
\makeatother
\subtitle{Provenance required by Hard Rule 7 for a deviation from Hard
Rule 2 (PROJECT spacetimes)}
\author{Prepared for Patrick Lee Nash, Ph.D.}
\date{November 12, 2026}

\begin{document}
\maketitle

\hypertarget{what-this-document-is-for}{%
\section{1. What this document is for}\label{what-this-document-is-for}}

Hard Rule 2 of the task specification says: \emph{push immediately to
the repository as soon as any object is updated}. Hard Rule 7 says that
an instruction which cannot be followed exactly must be explained in a
provenance. This document is that explanation. It records (i) what was
tried, with the exact error messages; (ii) why the pushes fail; (iii)
what the author decided; and (iv) the procedure by which every commit
nevertheless reaches \texttt{https://github.com/once-ere/spacetimes}
within about a minute of being made.

Nothing in this document concerns physics. It is the only document of
the project in which an instruction of the task specification is not
carried out literally by the assistant; the reason is a missing
permission, not a shortcut.

\hypertarget{the-three-places-where-git-can-run}{%
\section{2. The three places where git can
run}\label{the-three-places-where-git-can-run}}

Three different computers (or sandboxes) take part in this project. The
table fixes their names, because the rest of this document refers to
them.

\begin{longtable}[]{@{}
  >{\raggedright\arraybackslash}p{(\columnwidth - 4\tabcolsep) * \real{0.1765}}
  >{\raggedright\arraybackslash}p{(\columnwidth - 4\tabcolsep) * \real{0.3529}}
  >{\raggedright\arraybackslash}p{(\columnwidth - 4\tabcolsep) * \real{0.4706}}@{}}
\toprule\noalign{}
\begin{minipage}[b]{\linewidth}\raggedright
Name
\end{minipage} & \begin{minipage}[b]{\linewidth}\raggedright
What it is
\end{minipage} & \begin{minipage}[b]{\linewidth}\raggedright
What it can do
\end{minipage} \\
\midrule\noalign{}
\endhead
\bottomrule\noalign{}
\endlastfoot
\textbf{W}, the \emph{workspace} & A Linux machine (x86\_64) in the
cloud in which the assistant compiles and runs all code and builds all
documents. The repository clone is \texttt{/home/claude/spacetimes}. &
Clone and fetch from GitHub. \textbf{Cannot push} (Sec. 3). \\
\textbf{S}, the \emph{sandbox} & A Linux virtual machine (aarch64) on
the author's Mac. It sees the project folder
\texttt{/Users/nsh/Developer/github/spacetimes} and nothing else of the
Mac. & Run git inside the project folder. \textbf{Cannot push}: it holds
no GitHub credentials (Sec. 3). \\
\textbf{M}, the \emph{Mac} & The author's own macOS session (Terminal).
& \textbf{Can push}: it holds the author's GitHub credentials. \\
\end{longtable}

\hypertarget{what-was-tried-and-what-happened}{%
\section{3. What was tried, and what
happened}\label{what-was-tried-and-what-happened}}

All times are Pacific Standard Time on November 12, 2026, to within
about one minute.

\begin{longtable}[]{@{}
  >{\raggedright\arraybackslash}p{(\columnwidth - 6\tabcolsep) * \real{0.1463}}
  >{\raggedright\arraybackslash}p{(\columnwidth - 6\tabcolsep) * \real{0.1707}}
  >{\raggedright\arraybackslash}p{(\columnwidth - 6\tabcolsep) * \real{0.2195}}
  >{\raggedright\arraybackslash}p{(\columnwidth - 6\tabcolsep) * \real{0.4634}}@{}}
\toprule\noalign{}
\begin{minipage}[b]{\linewidth}\raggedright
Time
\end{minipage} & \begin{minipage}[b]{\linewidth}\raggedright
Where
\end{minipage} & \begin{minipage}[b]{\linewidth}\raggedright
Command
\end{minipage} & \begin{minipage}[b]{\linewidth}\raggedright
Result (verbatim)
\end{minipage} \\
\midrule\noalign{}
\endhead
\bottomrule\noalign{}
\endlastfoot
07:27 & S & \texttt{git\ push\ -\/-dry-run\ origin\ main} &
\texttt{fatal:\ could\ not\ read\ Username\ for\ \textquotesingle{}https://github.com\textquotesingle{}:\ terminal\ prompts\ disabled} \\
07:27 & W & attach repository with push access &
\texttt{pushes\ from\ this\ session\ will\ be\ refused:\ Claude\ doesn\textquotesingle{}t\ have\ GitHub\ access\ to\ once-ere/spacetimes\ for\ your\ organization} \\
07:35 & W & \texttt{git\ push\ -\/-dry-run\ origin\ HEAD:main} &
\texttt{remote:\ Claude\ doesn\textquotesingle{}t\ have\ GitHub\ access\ to\ once-ere/spacetimes\ for\ your\ organization.}
\ldots{}
\texttt{fatal:\ unable\ to\ access\ \textquotesingle{}https://github.com/once-ere/spacetimes/\textquotesingle{}:\ The\ requested\ URL\ returned\ error:\ 403} \\
07:41 & W & \texttt{git\ push\ -\/-dry-run\ origin\ HEAD:main} (after
the author tried to connect GitHub) & the same 403 refusal \\
07:41 & W & \texttt{gh\ auth\ status} &
\texttt{The\ token\ in\ GH\_TOKEN\ is\ invalid.} \\
07:41 & S & \texttt{git\ push\ -\/-dry-run\ origin\ main} &
\texttt{fatal:\ could\ not\ read\ Username\ for\ \textquotesingle{}https://github.com\textquotesingle{}:\ terminal\ prompts\ disabled} \\
\end{longtable}

\emph{Why W fails.} Pushes from the workspace go through the GitHub App
of the assistant's vendor. That App has not been granted access to
\texttt{once-ere/spacetimes}, and only the owner of the GitHub account
can grant it. Reading the repository works because the repository is
public.

\emph{Why S fails.} The sandbox is deliberately isolated from the rest
of the Mac. It has no access to the Mac's keychain, so git inside it has
no user name and no token.

\emph{Why M works.} The Mac's own git has working credentials for
\texttt{github.com/once-ere}: the remote-tracking log of another
repository of the author on the same machine records successful pushes
on November 9 and 10, 2026.

\hypertarget{the-authors-decisions}{%
\section{4. The author's decisions}\label{the-authors-decisions}}

\begin{longtable}[]{@{}
  >{\raggedright\arraybackslash}p{(\columnwidth - 4\tabcolsep) * \real{0.2308}}
  >{\raggedright\arraybackslash}p{(\columnwidth - 4\tabcolsep) * \real{0.3846}}
  >{\raggedright\arraybackslash}p{(\columnwidth - 4\tabcolsep) * \real{0.3846}}@{}}
\toprule\noalign{}
\begin{minipage}[b]{\linewidth}\raggedright
Time
\end{minipage} & \begin{minipage}[b]{\linewidth}\raggedright
Question
\end{minipage} & \begin{minipage}[b]{\linewidth}\raggedright
Decision
\end{minipage} \\
\midrule\noalign{}
\endhead
\bottomrule\noalign{}
\endlastfoot
07:36 & Pushes are refused; pause, or commit now and push later? &
``Commit now, push later.'' \\
07:41 & --- & ``I cannot get GitHub connected; try push.'' (Tried again;
refused again; Sec. 3.) \\
07:44 & --- & ``I will push from terminal; give me the exact
syntax.'' \\
\end{longtable}

\hypertarget{the-procedure-that-replaces-the-direct-push}{%
\section{5. The procedure that replaces the direct
push}\label{the-procedure-that-replaces-the-direct-push}}

Every update is committed at once in W. The commit is then carried to
the author's folder, where it becomes part of the local repository, and
M pushes it.

\textbf{Step 1 (W): commit.} Each update is committed in
\texttt{/home/claude/spacetimes} as soon as it is complete.

\textbf{Step 2 (W): bundle.} The new commits are written into one file,
a \emph{git bundle}, which is a complete and self-checking package of
commits:

\begin{Shaded}
\begin{Highlighting}[]
\FunctionTok{git}\NormalTok{ bundle create spacetimes\_update.bundle }\OperatorTok{\textless{}}\NormalTok{last{-}synchronized{-}commit}\OperatorTok{\textgreater{}}\NormalTok{..main}
\end{Highlighting}
\end{Shaded}

\textbf{Step 3 (W to the Mac): copy.} The bundle is copied to
\texttt{/Users/nsh/Developer/github/spacetimes/.git-transfer/spacetimes\_update.bundle}.
The folder \texttt{.git-transfer/} is excluded from version control by
the local file \texttt{.git/info/exclude}.

\textbf{Step 4 (S): fast-forward.} Inside the project folder the sandbox
verifies the bundle and advances the local branch to it. A
\emph{fast-forward} only adds commits; it never rewrites or discards
anything:

\begin{Shaded}
\begin{Highlighting}[]
\FunctionTok{git}\NormalTok{ bundle verify .git{-}transfer/spacetimes\_update.bundle}
\FunctionTok{git}\NormalTok{ fetch .git{-}transfer/spacetimes\_update.bundle main}
\FunctionTok{git}\NormalTok{ merge }\AttributeTok{{-}{-}ff{-}only}\NormalTok{ FETCH\_HEAD}
\end{Highlighting}
\end{Shaded}

After this step the files in the author's folder are exactly the files
of the commit, and \texttt{git\ rev-parse\ main} gives the same 40-digit
commit name in W and in the folder. This equality is checked after every
synchronization.

\textbf{Step 5 (M): push.} The author pushes from Terminal. Either once,

\begin{Shaded}
\begin{Highlighting}[]
\BuiltInTok{cd} \StringTok{"/Users/nsh/Developer/github/spacetimes"}
\FunctionTok{git}\NormalTok{ push origin main}
\end{Highlighting}
\end{Shaded}

or continuously, with a loop that tries a push every 60 seconds and is
silent unless there is an error (stop it with Control-C):

\begin{Shaded}
\begin{Highlighting}[]
\BuiltInTok{cd} \StringTok{"/Users/nsh/Developer/github/spacetimes"}
\ControlFlowTok{while} \FunctionTok{true}\KeywordTok{;} \ControlFlowTok{do} \FunctionTok{git}\NormalTok{ push }\AttributeTok{{-}q}\NormalTok{ origin main}\KeywordTok{;} \FunctionTok{sleep}\NormalTok{ 60}\KeywordTok{;} \ControlFlowTok{done}
\end{Highlighting}
\end{Shaded}

With the loop running, every commit reaches GitHub within about one
minute of Step 4, which is as close to ``push immediately'' as the
missing permission allows.

\hypertarget{how-to-check-that-nothing-is-waiting}{%
\section{6. How to check that nothing is
waiting}\label{how-to-check-that-nothing-is-waiting}}

\begin{Shaded}
\begin{Highlighting}[]
\BuiltInTok{cd} \StringTok{"/Users/nsh/Developer/github/spacetimes"}
\FunctionTok{git}\NormalTok{ status }\AttributeTok{{-}sb}
\end{Highlighting}
\end{Shaded}

The first line of the answer is \texttt{\#\#\ main...origin/main} when
everything has been pushed, and
\texttt{\#\#\ main...origin/main\ {[}ahead\ N{]}} when \texttt{N}
commits are still waiting.

\hypertarget{what-is-deliberately-not-in-the-repository}{%
\section{7. What is deliberately not in the
repository}\label{what-is-deliberately-not-in-the-repository}}

The files that were already in the folder when the project started are
not committed:

\begin{itemize}
\tightlist
\item
  the author's Mathematica notebooks (\texttt{4+4.nb}, \texttt{7+1.nb}
  and the others);
\item
  the nine PDF files of published journal articles, which are not
  redistributed here because the right to republish them rests with
  their publishers;
\item
  the author's task specification and working notes
  (\texttt{prompt*.txt}, \texttt{notes\_for\_prompt*.txt}), which the
  author's \texttt{.gitignore} declares private.
\end{itemize}

They remain in the folder as untracked files. Wherever the project
depends on one of them, the provenance documents identify it by its
SHA-256 checksum, so that the exact input can be recognized without
being published. If the author wants the notebooks in the repository,
one command adds them: \texttt{git\ add\ 4+4.nb\ 7+1.nb}, followed by a
commit.

\hypertarget{if-the-permission-is-granted-later}{%
\section{8. If the permission is granted
later}\label{if-the-permission-is-granted-later}}

If the GitHub App is given access to \texttt{once-ere/spacetimes}, the
workspace W can push directly. Steps 2 to 5 then become unnecessary; the
local folder is brought up to date with \texttt{git\ pull\ -\/-ff-only},
and this document will be amended to say so.

\end{document}
